
\documentclass[aps,preprint]{revtex4}
%%%%%%%%%%%%%%%%%%%%%%%%%%%%%%%%%%%%%%%%%%%%%%%%%%%%%%%%%%%%%%%%%%%%%%%%%%%%%%%%%%%%%%%%%%%%%%%%%%%%%%%%%%%%%%%%%%%%%%%%%%%%
\usepackage{amsfonts}
\usepackage{amsmath}

\setcounter{MaxMatrixCols}{10}
%TCIDATA{OutputFilter=LATEX.DLL}
%TCIDATA{Version=4.10.0.2363}
%TCIDATA{Created=Thursday, January 03, 2002 09:02:25}
%TCIDATA{LastRevised=Saturday, June 21, 2003 12:00:21}
%TCIDATA{<META NAME="GraphicsSave" CONTENT="32">}
%TCIDATA{<META NAME="DocumentShell" CONTENT="Articles\SW\REVTeX 4">}
%TCIDATA{Language=American English}
%TCIDATA{CSTFile=revtex4.cst}

\newtheorem{theorem}{Theorem}
\newtheorem{acknowledgement}[theorem]{Acknowledgement}
\newtheorem{algorithm}[theorem]{Algorithm}
\newtheorem{axiom}[theorem]{Axiom}
\newtheorem{claim}[theorem]{Claim}
\newtheorem{conclusion}[theorem]{Conclusion}
\newtheorem{condition}[theorem]{Condition}
\newtheorem{conjecture}[theorem]{Conjecture}
\newtheorem{corollary}[theorem]{Corollary}
\newtheorem{criterion}[theorem]{Criterion}
\newtheorem{definition}[theorem]{Definition}
\newtheorem{example}[theorem]{Example}
\newtheorem{exercise}[theorem]{Exercise}
\newtheorem{lemma}[theorem]{Lemma}
\newtheorem{notation}[theorem]{Notation}
\newtheorem{problem}[theorem]{Problem}
\newtheorem{proposition}[theorem]{Proposition}
\newtheorem{remark}[theorem]{Remark}
\newtheorem{solution}[theorem]{Solution}
\newtheorem{summary}[theorem]{Summary}
\newenvironment{proof}[1][Proof]{\noindent\textbf{#1.} }{\ \rule{0.5em}{0.5em}}
\input{tcilatex}

\begin{document}

\preprint{}
\title[Geminal Density Operators]{Geminal Density Operators}
\author{Brian Weiner}
\affiliation{PSU}
\date{06/17/03}
\startpage{1}
\endpage{2}
\maketitle
\tableofcontents

\section{Two Particle Density Operators\label{Two Particle}}

\subsection{Orthonormal Basis}

\subsubsection{Contraction}

The normalized geminal $\left\vert g\right\rangle $ can be expanded on an
unordered orthonormal basis as 
\begin{equation}
\left\vert g\right\rangle =\frac{1}{2}\sum_{1\leq j_{1},j_{2}\leq
r}c_{j_{1}j_{2}}\left\vert \psi _{j_{1}}\otimes \psi _{j_{2}}\right\rangle
;\quad c_{j_{1}j_{2}}=-c_{j_{2}j_{1}}
\end{equation}%
and its density operator as 
\begin{equation}
\left\vert g\right\rangle \left\langle g\right\vert =\frac{1}{4}\sum
_{\substack{ 1\leq j_{1},j_{2}\leq r \\ 1\leq k_{1},k_{2}\leq r}}%
c_{j_{1}j_{2}}\left\vert \psi _{j_{1}}\otimes \psi _{j_{2}}\right\rangle
\left\langle \psi _{k_{1}}\otimes \psi _{k_{2}}\right\vert
c_{k_{1}k_{2}}^{\ast }
\end{equation}%
As $\left\langle g|g\right\rangle =1$ one observes that 
\begin{equation}
\func{Tr}\left\{ \mathbf{cc}^{\dagger }\right\} =4
\end{equation}%
The expansion in terms of the tensor product basis can be compared to 
\begin{equation}
\left\vert g\right\rangle =\frac{1}{2}\sum_{1\leq j_{1},j_{2}\leq r}\tilde{c}%
_{j_{1}j_{2}}\left\vert \psi _{j_{1}}\wedge \psi _{j_{2}}\right\rangle
;\quad \tilde{c}_{j_{1}j_{2}}=-\tilde{c}_{j_{2}j_{1}}
\end{equation}%
in terms of the antisymmetrized tensor product basis 
\begin{equation}
\left\vert \psi _{j_{1}}\wedge \psi _{j_{2}}\right\rangle =\frac{1}{\sqrt{2}}%
\left( \left\vert \psi _{j_{1}}\otimes \psi _{j_{2}}\right\rangle
-\left\vert \psi _{j_{2}}\otimes \psi _{j_{1}}\right\rangle \right) 
\end{equation}%
Evaluating the antisymmetric expansion we obtain 
\begin{eqnarray}
\frac{1}{2}\sum_{1\leq j_{1},j_{2}\leq r}\tilde{c}_{j_{1}j_{2}}\left\vert
\psi _{j_{1}}\wedge \psi _{j_{2}}\right\rangle  &=&\frac{1}{2}\sum_{1\leq
j_{1},j_{2}\leq r}\tilde{c}_{j_{1}j_{2}}\frac{1}{\sqrt{2}}\left( \left\vert
\psi _{j_{1}}\otimes \psi _{j_{2}}\right\rangle -\left\vert \psi
_{j_{2}}\otimes \psi _{j_{1}}\right\rangle \right)   \notag \\
&=&\frac{1}{2}\frac{1}{\sqrt{2}}\sum_{1\leq j_{1},j_{2}\leq r}\tilde{c}%
_{j_{1}j_{2}}\left\vert \psi _{j_{1}}\otimes \psi _{j_{2}}\right\rangle -%
\frac{1}{2}\frac{1}{\sqrt{2}}\sum_{1\leq j_{1},j_{2}\leq r}\tilde{c}%
_{j_{1}j_{2}}\left\vert \psi _{j_{2}}\otimes \psi _{j_{1}}\right\rangle  
\notag \\
&=&\frac{1}{2}\frac{1}{\sqrt{2}}\sum_{1\leq j_{1},j_{2}\leq r}\tilde{c}%
_{j_{1}j_{2}}\left\vert \psi _{j_{1}}\otimes \psi _{j_{2}}\right\rangle +%
\frac{1}{2}\frac{1}{\sqrt{2}}\sum_{1\leq j_{1},j_{2}\leq r}\tilde{c}%
_{j_{1}j_{2}}\left\vert \psi _{j_{1}}\otimes \psi _{j_{2}}\right\rangle  
\notag \\
&=&\frac{2}{\sqrt{2}}\sum_{1\leq j_{1},j_{2}\leq r}\tilde{c}%
_{j_{1}j_{2}}\left\vert \psi _{j_{1}}\otimes \psi _{j_{2}}\right\rangle 
\end{eqnarray}%
which shows that 
\begin{equation}
c_{j_{1}j_{2}}=\sqrt{2}\tilde{c}_{j_{1}j_{2}}
\end{equation}%
The contraction map \textrm{L}$_{2}^{1}$ acting on $\left\vert
g\right\rangle \left\langle g\right\vert $ produces 
\begin{eqnarray}
\mathrm{L}_{2}^{1}\left( \left\vert g\right\rangle \left\langle g\right\vert
\right)  &=&\frac{1}{4}\sum_{\substack{ 1\leq j_{1},j_{2}\leq r \\ 1\leq
k_{1},k_{2}\leq r}}c_{j_{1}j_{2}}c_{k_{1}k_{2}}^{\ast }\left\langle \psi
_{k_{2}}|\psi _{j_{2}}\right\rangle \left\vert \psi _{j_{1}}\right\rangle
\left\langle \psi _{k_{1}}\right\vert =\frac{1}{4}\sum_{\substack{ 1\leq
j_{1},j_{2}\leq r \\ 1\leq k_{1},k_{2}\leq r}}c_{j_{1}j_{2}}c_{k_{1}k_{2}}^{%
\ast }\delta _{k_{2}j_{2}}\left\vert \psi _{j_{1}}\right\rangle \left\langle
\psi _{k_{1}}\right\vert   \notag \\
&=&\frac{1}{4}\sum_{\substack{ 1\leq j_{1},j_{2}\leq r \\ 1\leq k_{1}\leq r}}%
c_{j_{1}j_{2}}c_{k_{1}j_{2}}^{\ast }\left\vert \psi _{j_{1}}\right\rangle
\left\langle \psi _{k_{1}}\right\vert =\frac{1}{4}\sum_{1\leq
j_{1},j_{2}\leq r}\left( \mathbf{cc}^{\dagger }\right)
_{j_{1}k_{1}}\left\vert \psi _{j_{1}}\right\rangle \left\langle \psi
_{k_{1}}\right\vert 
\end{eqnarray}%
The First Order Reduced Density Operator (FORDO) $D^{1}\left( g\right) $
normalized to the number of electrons in the system is given by 
\begin{equation}
D^{1}\left( g\right) =2\mathrm{L}_{2}^{1}\left( \left\vert g\right\rangle
\left\langle g\right\vert \right) 
\end{equation}%
and the First Order Reduced Density Matrix (FORDM), $\mathbb{D}_{\psi
}^{1}\left( g\right) ,$ which is the \emph{representation} of $D^{1}\left(
g\right) $ in the $\mathbf{\psi }$ basis by 
\begin{equation}
\mathbb{D}_{\psi }^{1}\left( g\right) =\frac{1}{2}\mathbf{cc}^{\dagger }=%
\mathbf{\tilde{c}\tilde{c}}^{\dagger }
\end{equation}%
both of which have trace equal to $2$ i.e. 
\begin{equation}
\func{Tr}\left\{ D^{1}\left( g\right) \right\} =\frac{1}{2}\func{Tr}\left\{ 
\mathbf{cc}^{\dagger }\right\} =\func{Tr}\left\{ \mathbf{\tilde{c}\tilde{c}}%
^{\dagger }\right\} =2
\end{equation}%
One can also observe that 
\begin{equation}
\mathbb{D}_{\psi }^{1}\left( g\right) ^{\ast }=\frac{1}{2}\mathbf{c}%
^{\dagger }\mathbf{c}
\end{equation}%
as $\mathbf{c}$ is antisymmetric 
\begin{equation}
\mathbb{D}_{\psi }^{1}\left( g\right) =\frac{1}{2}\mathbf{cc}^{\dagger }=-%
\frac{1}{2}\mathbf{cc}^{\ast }=\frac{1}{2}\mathbf{c}^{t}\mathbf{c}^{\ast }
\end{equation}%
The NGSO's $\left\{ \mathbf{w}_{j};1\leq j\leq r\right\} $ of $\mathbb{D}%
_{\psi }^{1}\left( g\right) $ are given by 
\begin{equation}
\mathbb{D}_{\psi }^{1}\left( g\right) \mathbf{w}_{j}=n_{j}\mathbf{w}_{j}
\end{equation}%
and 
\begin{equation}
\mathbb{W}^{\dagger }\mathbb{D}_{\psi }^{1}\left( g\right) \mathbb{W=}\left( 
\begin{array}{cc}
\mathbf{n} & \mathbf{0} \\ 
\mathbf{0} & \mathbf{n}%
\end{array}%
\right) 
\end{equation}%
where 
\begin{equation}
\mathbf{n=}\left( 
\begin{array}{ccc}
n_{1} & 0 & 0 \\ 
0 & \ddots  & 0 \\ 
0 & 0 & n_{s}%
\end{array}%
\right) ;\qquad n_{j}\geq 0,\quad r=2s
\end{equation}

\subsection{Non Orthonormal Basis}

\subsubsection{Basis Transformation}

If the basis $\left\vert \mathbf{\psi }\right\rangle $ is a transformation
of another basis $\left\vert \mathbf{\varphi }\right\rangle $ that is not
necessarily orthonormal then.%
\begin{equation}
\left\vert \mathbf{\psi }\right\rangle =V\left\vert \mathbf{\varphi }%
\right\rangle =\left\vert \mathbf{\varphi }\right\rangle \mathbb{V}_{\varphi
}
\end{equation}%
i.e. 
\begin{equation}
\left\vert \psi _{j}\right\rangle =V\left\vert \varphi _{j}\right\rangle
=\sum_{1\leq k\leq r}\left\vert \varphi _{k}\right\rangle V_{kj};\qquad
1\leq j\leq r
\end{equation}%
We also have the following relationships 
\begin{eqnarray}
\left\langle \mathbf{\psi }|\mathbf{\psi }\right\rangle &=&\mathbb{I}%
_{r}=\left\langle \mathbf{\varphi }|V^{\dagger }V\mathbf{\varphi }%
\right\rangle =\mathbb{V}_{\varphi }^{\dagger }\left\langle \mathbf{\varphi }%
|\mathbf{\varphi }\right\rangle \mathbb{V}_{\varphi }\equiv \mathbb{V}%
_{\varphi }^{\dagger }\mathbb{SV}_{\varphi }  \notag \\
\mathbb{V}_{\varphi }^{-1\dagger }\mathbb{V}_{\varphi }^{-1} &=&\mathbb{S} 
\notag \\
\mathbb{S}^{-1} &=&\mathbb{V}_{\varphi }\mathbb{V}_{\varphi }^{\dagger } 
\notag \\
\mathbb{V}_{\varphi }^{-1} &=&\mathbb{V}_{\varphi }^{\dagger }\mathbb{S} 
\notag \\
\mathbb{V}_{\varphi }^{\dagger } &=&\mathbb{S}^{-1}\mathbb{V}_{\varphi }^{-1}
\end{eqnarray}%
The geminal can be expressed on the general basis as 
\begin{eqnarray}
\left\vert g\right\rangle &=&\frac{1}{2}\sum_{1\leq j_{1},j_{2}\leq
r}c_{j_{1}j_{2}}\left\vert V\varphi _{j_{1}}V\varphi _{j_{2}}\right\rangle =%
\frac{1}{2}\sum_{1\leq k_{1},k_{2}\leq r}\gamma _{k_{1}k_{2}}\left\vert
\varphi _{k_{1}}\varphi _{k_{2}}\right\rangle  \notag \\
&=&\frac{1}{2}\sum_{\substack{ 1\leq j_{1},j_{2}\leq r  \\ 1\leq
k_{1},k_{2}\leq r}}c_{j_{1}j_{2}}V_{k_{1}j_{1}}V_{k_{2}j_{2}}\left\vert
\varphi _{k_{1}}\varphi _{k_{2}}\right\rangle =\frac{1}{2}\sum_{\substack{ %
1\leq j_{1},j_{2}\leq r  \\ 1\leq k_{1},k_{2}\leq r}}%
V_{k_{1}j_{1}}c_{j_{1}j_{2}}V_{j_{2}k_{2}}^{t}\left\vert \varphi
_{k_{1}}\varphi _{k_{2}}\right\rangle  \notag \\
&=&\frac{1}{2}\sum_{1\leq k_{1},k_{2}\leq r}\left( \mathbb{V}_{\varphi }%
\mathbf{c}\mathbb{V}_{\varphi }^{t}\right) _{k_{1}k_{2}}\left\vert \varphi
_{k_{1}}\varphi _{k_{2}}\right\rangle
\end{eqnarray}%
i.e. 
\begin{equation}
\mathbf{\gamma =}\mathbb{V}_{\varphi }\mathbf{c}\mathbb{V}_{\varphi }^{t}=-%
\mathbf{\gamma }^{t}
\end{equation}%
We can also note that 
\begin{eqnarray}
\left\langle \mathbf{\varphi \otimes \varphi }|g\right\rangle &=&\frac{1}{2}%
\mathbb{S}\mathbf{\gamma }\mathbb{S}^{t}=\frac{1}{2}\mathbb{SV}_{\varphi }%
\mathbf{c}\mathbb{V}_{\varphi }^{t}\mathbb{S}^{t}  \notag \\
\left\langle \mathbf{\psi \otimes \psi }|g\right\rangle &=&\frac{1}{2}%
\mathbf{c}
\end{eqnarray}

\subsubsection{Contraction}

On the general basis 
\begin{eqnarray}
\mathrm{L}_{2}^{1}\left( \left\vert g\right\rangle \left\langle g\right\vert
\right) &=&\frac{1}{4}\sum_{\substack{ 1\leq j_{1},j_{2}\leq r  \\ 1\leq
k_{1},k_{2}\leq r}}\gamma _{j_{1}j_{2}}\gamma _{k_{1}k_{2}}^{\ast
}\left\langle \varphi _{k_{2}}|\varphi _{j_{2}}\right\rangle \left\vert
\varphi _{j_{1}}\right\rangle \left\langle \varphi _{k_{1}}\right\vert =%
\frac{1}{4}\sum_{\substack{ 1\leq j_{1},j_{2}\leq r  \\ 1\leq
k_{1},k_{2}\leq r}}\gamma _{j_{1}j_{2}}\gamma _{k_{1}k_{2}}^{\ast
}S_{k_{2}j_{2}}\left\vert \varphi _{j_{1}}\right\rangle \left\langle \varphi
_{k_{1}}\right\vert  \notag \\
&=&\frac{1}{4}\sum_{1\leq j_{1},j_{2}\leq r}\left( \mathbf{\gamma }\mathbb{S}%
^{t}\mathbf{\gamma }^{\dagger }\right) _{j_{1}k_{1}}\left\vert \varphi
_{j_{1}}\right\rangle \left\langle \varphi _{k_{1}}\right\vert =\frac{1}{4}%
\sum_{1\leq j_{1},j_{2}\leq r}\left( \mathbf{\gamma }^{t}\mathbb{S}^{\ast }%
\mathbf{\gamma }^{\ast }\right) _{j_{1}k_{1}}\left\vert \varphi
_{j_{1}}\right\rangle \left\langle \varphi _{k_{1}}\right\vert
\end{eqnarray}%
Thus the (FORDM) in the general basis is (see section $\left( \ref{Change}%
\right) $) 
\begin{equation}
\mathbb{D}_{\varphi }^{1}\left( g\right) =\frac{1}{2}\mathbf{\gamma \,}%
\mathbb{S}^{t}\mathbf{\gamma }^{\dagger }\mathbb{S}
\end{equation}%
Hence we have 
\begin{eqnarray}
\mathbb{D}_{\varphi }^{1}\left( g\right) &=&\mathbb{S}^{-1}\left\langle 
\mathbf{\varphi }|D^{1}\left( g\right) \mathbf{\varphi }\right\rangle  \notag
\\
\mathbb{D}_{\psi }^{1}\left( g\right) &=&\left\langle \mathbf{\psi }%
|D^{1}\left( g\right) \mathbf{\psi }\right\rangle  \notag \\
\mathbb{D}_{\varphi }^{1}\left( g\right) &=&\mathbb{V}_{\varphi }\mathbb{D}%
_{\psi }^{1}\left( g\right) \mathbb{V}_{\varphi }^{-1}  \label{dentran}
\end{eqnarray}%
and 
\begin{eqnarray}
\left\langle \mathbf{\psi }|D^{1}\left( g\right) \mathbf{\psi }\right\rangle
&=&\frac{1}{2}\mathbf{cc}^{\dagger }=\mathbb{D}_{\psi }^{1}\left( g\right) 
\notag \\
\left\langle \mathbf{\varphi }|D^{1}\left( g\right) \mathbf{\varphi }%
\right\rangle &=&\frac{1}{2}\mathbb{S}\mathbf{\gamma \,}\mathbb{S}^{t}%
\mathbf{\gamma }^{\dagger }\mathbb{S}\neq \mathbb{D}_{\varphi }^{1}\left(
g\right)
\end{eqnarray}%
as (see section $\left( \ref{Change}\right) $) 
\begin{equation}
\mathbb{D}_{\varphi }^{1}\left( g\right) =\mathbb{S}^{-1}\left\langle 
\mathbf{\varphi }|D^{1}\left( g\right) \mathbf{\varphi }\right\rangle =\frac{%
1}{2}\mathbf{\gamma \,}\mathbb{S}^{t}\mathbf{\gamma }^{\dagger }\mathbb{S}
\end{equation}%
where 
\begin{equation}
\mathbf{\gamma }=\mathbb{V}_{\varphi }\mathbf{c}\mathbb{V}_{\varphi }^{t}
\end{equation}%
as test of consistency we see that 
\begin{equation}
\mathbb{V}_{\varphi }^{-1}\mathbb{D}_{\varphi }^{1}\left( g\right) \mathbb{V}%
_{\varphi }=\frac{1}{2}\mathbb{V}_{\varphi }^{-1}\mathbf{\gamma \,}\mathbb{S}%
^{t}\mathbf{\gamma }^{\dagger }\mathbb{SV}_{\varphi }=\frac{1}{2}\mathbb{V}%
_{\varphi }^{-1}\mathbb{V}_{\varphi }\mathbf{c}\mathbb{V}_{\varphi }^{t}%
\mathbf{\,}\mathbb{S}^{t}\mathbb{V}_{\varphi }^{\ast }\mathbf{c}^{\dagger }%
\mathbb{V}_{\varphi }^{\dagger }\mathbb{SV}_{\varphi }=\frac{1}{2}\mathbf{cc}%
^{\dagger }=\mathbb{D}_{\psi }^{1}\left( g\right)
\end{equation}%
as it should.

The NGSO's $\left\{ \mathbf{w}_{\varphi j};1\leq j\leq r\right\} $ of $%
\mathbb{D}_{\varphi }^{1}\left( g\right) $ are given by 
\begin{equation}
\mathbb{D}_{\varphi }^{1}\left( g\right) \mathbf{w}_{\varphi j}=n_{j}\mathbf{%
w}_{\varphi j}
\end{equation}%
where 
\begin{equation}
\mathbf{W}_{\varphi }^{\dagger }\mathbb{S}\mathbf{W}_{\varphi }=\mathbb{I}%
_{r}
\end{equation}%
$\lceil \left\langle w_{j}|w_{k}\right\rangle =\mathbf{w}_{\varphi
j}^{\dagger }\mathbb{S}\mathbf{w}_{\varphi k}\rfloor $ hence%
\begin{equation}
\mathbf{W}_{\varphi }^{\dagger }\mathbb{SD}_{\varphi }^{1}\left( g\right) 
\mathbf{W}_{\varphi }=\left( 
\begin{array}{cc}
\mathbf{n} & \mathbf{0} \\ 
\mathbf{0} & \mathbf{n}%
\end{array}%
\right)
\end{equation}

\section{Special Orthonormal Basis\label{Special}}

We can choose the orthonormal basis to be the symmetrical orthonormalized
basis i.e.$\mathbb{V}_{\varphi }=\mathbb{S}^{-\frac{1}{2}}$ so that 
\begin{equation}
\left\vert \mathbf{\psi }\right\rangle =S^{-\frac{1}{2}}\left\vert \mathbf{%
\varphi }\right\rangle =\left\vert \mathbf{\varphi }\right\rangle \mathbb{S}%
^{-\frac{1}{2}}
\end{equation}%
then 
\begin{equation}
\mathbb{D}_{\psi }^{1}\left( g\right) =\mathbb{S}^{\frac{1}{2}}\mathbb{D}%
_{\varphi }^{1}\left( g\right) \mathbb{S}^{-\frac{1}{2}}=\mathbb{S}^{\frac{1%
}{2}}\frac{1}{2}\mathbf{\gamma \,}\mathbb{S}^{t}\mathbf{\gamma }^{\dagger }%
\mathbb{SS}^{-\frac{1}{2}}=\frac{1}{2}\mathbb{S}^{\frac{1}{2}}\mathbf{\gamma
\,}\mathbb{S}^{t}\mathbf{\gamma }^{\dagger }\mathbb{S}^{\frac{1}{2}}
\end{equation}%
and solve the eigenvalue problem in this basis to give 
\begin{equation}
\mathbf{W}_{\psi }^{\dagger }\mathbb{D}_{\psi }^{1}\left( g\right) \mathbf{W}%
_{\psi }=\frac{1}{2}\mathbf{W}_{\psi }^{\dagger }\mathbb{S}^{\frac{1}{2}}%
\mathbf{\gamma \,}\mathbb{S}^{t}\mathbf{\gamma }^{\dagger }\mathbb{S}^{\frac{%
1}{2}}\mathbf{W}_{\psi }=\left( 
\begin{array}{cc}
\mathbf{n} & \mathbf{0} \\ 
\mathbf{0} & \mathbf{n}%
\end{array}%
\right)
\end{equation}%
where 
\begin{equation}
\mathbf{W}_{\psi }^{\dagger }\mathbf{W}_{\psi }=\mathbf{W}_{\psi }\mathbf{W}%
_{\psi }^{\dagger }=\mathbb{I}_{r}
\end{equation}%
The geminal coefficients in this orthogonal basis are thus 
\begin{equation}
\mathbf{c}=\mathbb{S}^{\frac{1}{2}}\mathbf{\gamma }\mathbb{S}^{\frac{1}{2}t}
\end{equation}%
giving 
\begin{equation}
\mathbf{cc}^{\dagger }=\mathbb{S}^{\frac{1}{2}}\mathbf{\gamma }\mathbb{S}^{%
\frac{1}{2}t}\mathbb{S}^{\frac{1}{2}t}\mathbf{\gamma }^{\dagger }\mathbb{S}^{%
\frac{1}{2}}=\mathbb{S}^{\frac{1}{2}}\mathbf{\gamma }\mathbb{S}^{t}\mathbf{%
\gamma }^{\dagger }\mathbb{S}^{\frac{1}{2}}
\end{equation}%
which is consistent with eq. $\left( \ref{dentran}\right) $

One can further observe that 
\begin{eqnarray}
\mathbb{D}_{\psi }^{1}\left( g\right) \mathbf{W}_{\psi } &=&\mathbf{W}_{\psi
}\left( 
\begin{array}{cc}
\mathbf{n} & \mathbf{0} \\ 
\mathbf{0} & \mathbf{n}%
\end{array}%
\right) =  \notag \\
\frac{1}{2}\mathbb{S}^{\frac{1}{2}}\mathbf{\gamma \,}\mathbb{S}^{t}\mathbf{%
\gamma }^{\dagger }\mathbb{S}^{\frac{1}{2}}\mathbf{W}_{\psi } &=&\mathbf{W}%
_{\psi }\left( 
\begin{array}{cc}
\mathbf{n} & \mathbf{0} \\ 
\mathbf{0} & \mathbf{n}%
\end{array}%
\right) \Rightarrow  \notag \\
\frac{1}{2}\mathbb{S}^{\frac{1}{2}}\mathbf{\gamma \,}\mathbb{S}^{t}\mathbf{%
\gamma }^{\dagger }\mathbb{SS}^{-\frac{1}{2}}\mathbf{W}_{\psi } &=&\mathbf{W}%
_{\psi }\left( 
\begin{array}{cc}
\mathbf{n} & \mathbf{0} \\ 
\mathbf{0} & \mathbf{n}%
\end{array}%
\right) \Rightarrow  \notag \\
\frac{1}{2}\mathbf{\gamma \,}\mathbb{S}^{t}\mathbf{\gamma }^{\dagger }%
\mathbb{SS}^{-\frac{1}{2}}\mathbf{W}_{\psi } &=&\mathbb{S}^{-\frac{1}{2}}%
\mathbf{W}_{\psi }\left( 
\begin{array}{cc}
\mathbf{n} & \mathbf{0} \\ 
\mathbf{0} & \mathbf{n}%
\end{array}%
\right) =  \notag \\
\frac{1}{2}\mathbf{\gamma \,}\mathbb{S}^{t}\mathbf{\gamma }^{\dagger }%
\mathbb{S}\mathbf{W}_{\varphi } &=&\mathbf{W}_{\varphi }\left( 
\begin{array}{cc}
\mathbf{n} & \mathbf{0} \\ 
\mathbf{0} & \mathbf{n}%
\end{array}%
\right) =  \notag \\
\mathbb{D}_{\varphi }^{1}\left( g\right) \mathbf{W}_{\varphi } &=&\mathbf{W}%
_{\varphi }\left( 
\begin{array}{cc}
\mathbf{n} & \mathbf{0} \\ 
\mathbf{0} & \mathbf{n}%
\end{array}%
\right)
\end{eqnarray}%
where 
\begin{eqnarray}
\mathbf{W}_{\varphi } &=&\mathbb{S}^{-\frac{1}{2}}\mathbf{W}_{\psi }  \notag
\\
\mathbb{D}_{\varphi }^{1}\left( g\right) &=&\frac{1}{2}\mathbf{\gamma \,}%
\mathbb{S}^{t}\mathbf{\gamma }^{\dagger }\mathbb{S}
\end{eqnarray}%
and is consistent with 
\begin{equation}
\mathbf{W}_{\varphi }^{\dagger }\mathbb{SD}_{\varphi }^{1}\left( g\right) 
\mathbf{W}_{\varphi }=\mathbf{W}_{\varphi }^{\dagger }\mathbb{S}\mathbf{W}%
_{\varphi }\left( 
\begin{array}{cc}
\mathbf{n} & \mathbf{0} \\ 
\mathbf{0} & \mathbf{n}%
\end{array}%
\right) =\left( 
\begin{array}{cc}
\mathbf{n} & \mathbf{0} \\ 
\mathbf{0} & \mathbf{n}%
\end{array}%
\right)
\end{equation}%
If $\mathbf{W}_{\psi }$ are \emph{the} canonical orbitals of $g$ then 
\begin{eqnarray}
&&\mathbf{W}_{\psi }^{\dagger }\mathbf{cW}_{\psi }^{\ast }=\left( 
\begin{array}{cc}
\mathbf{0} & \mathbf{n}^{\frac{1}{2}} \\ 
\mathbf{-n}^{\frac{1}{2}} & \mathbf{0}%
\end{array}%
\right)  \notag \\
&\equiv &\mathbf{W}_{\psi }^{\dagger }\mathbb{S}^{\frac{1}{2}}\mathbf{\gamma 
}\mathbb{S}^{\frac{1}{2}t}\mathbf{W}_{\psi }^{\ast }=\left( 
\begin{array}{cc}
\mathbf{0} & \mathbf{n}^{\frac{1}{2}} \\ 
\mathbf{-n}^{\frac{1}{2}} & \mathbf{0}%
\end{array}%
\right)  \notag \\
&\equiv &\mathbf{W}_{\varphi }^{\dagger }\mathbb{S}\mathbf{\gamma }\mathbb{S}%
^{t}\mathbf{W}_{\varphi }^{\ast }=\left( 
\begin{array}{cc}
\mathbf{0} & \mathbf{n}^{\frac{1}{2}} \\ 
\mathbf{-n}^{\frac{1}{2}} & \mathbf{0}%
\end{array}%
\right)
\end{eqnarray}%
Noting that 
\begin{eqnarray}
\mathbf{W}_{\varphi }^{\dagger }\mathbb{S}\mathbf{W}_{\varphi } &=&\mathbb{I}
\notag \\
\mathbf{W}_{\varphi }^{-1} &=&\mathbf{W}_{\varphi }^{\dagger }\mathbb{S} 
\notag \\
\mathbf{W}_{\varphi }^{\ast -1} &=&\mathbf{W}_{\varphi }^{t}\mathbb{S}^{t} 
\notag \\
\mathbf{W}_{\varphi }^{\dagger -1} &=&\mathbb{S}\mathbf{W}_{\varphi }
\end{eqnarray}%
we obtain 
\begin{equation}
\mathbf{c}=\mathbf{W}_{\psi }\left( 
\begin{array}{cc}
\mathbf{0} & \mathbf{n}^{\frac{1}{2}} \\ 
\mathbf{-n}^{\frac{1}{2}} & \mathbf{0}%
\end{array}%
\right) \mathbf{W}_{\psi }^{t}
\end{equation}%
and 
\begin{equation}
\mathbf{\gamma }=\mathbf{W}_{\varphi }\left( 
\begin{array}{cc}
\mathbf{0} & \mathbf{n}^{\frac{1}{2}} \\ 
\mathbf{-n}^{\frac{1}{2}} & \mathbf{0}%
\end{array}%
\right) \mathbf{W}_{\varphi }^{t}
\end{equation}

\section{Change of Basis\label{Change}}

\subsubsection{Basis transformations}

The transformation between the $\mathbf{\varphi }$ and $\mathbf{\psi }$
basis induce the following transformations between the representations: 
\begin{equation}
A\left\vert \mathbf{\psi }\right\rangle =\left\vert \mathbf{\psi }%
\right\rangle \mathbb{A}_{\psi }=\left\vert \mathbf{\varphi }\right\rangle 
\mathbb{V}_{\varphi }\mathbb{A}_{\psi }=A\left\vert \mathbf{\varphi }%
\right\rangle \mathbb{V}_{\varphi }=\left\vert \mathbf{\varphi }%
\right\rangle \mathbb{A}_{\varphi }\mathbb{V}_{\varphi }
\end{equation}%
thus 
\begin{equation}
\mathbb{V}_{\varphi }\mathbb{A}_{\psi }=\mathbb{A}_{\varphi }\mathbb{V}%
_{\varphi }\Rightarrow \mathbb{A}_{\varphi }=\mathbb{V}_{\varphi }\mathbb{A}%
_{\psi }\mathbb{V}_{\varphi }^{-1}\text{ }\Longleftrightarrow \text{ }%
\mathbb{A}_{\psi }=\mathbb{V}_{\varphi }^{-1}\mathbb{A}_{\varphi }\mathbb{V}%
_{\varphi }
\end{equation}%
where $\mathbb{A}_{\varphi }$ and $\mathbb{A}_{\psi }$ are the
representations of the operator on the basis $\mathbf{\varphi }$ and $%
\mathbf{\psi }$ respectively. The relationship between the matrix of an
operator and its representation can be easily obtained by%
\begin{equation}
\left\langle \mathbf{\varphi }|A\mathbf{\varphi }\right\rangle =\left\langle 
\mathbf{\varphi }|\mathbf{\varphi }\right\rangle \mathbb{A}_{\varphi }\equiv 
\mathbb{SA}_{\varphi }\Rightarrow \mathbb{A}_{\varphi }=\mathbb{S}%
^{-1}\left\langle \mathbf{\varphi }|A\mathbf{\varphi }\right\rangle
\end{equation}%
and 
\begin{equation}
\left\langle \mathbf{\psi }|A\mathbf{\psi }\right\rangle =\left\langle 
\mathbf{\psi }|\mathbf{\psi }\right\rangle \mathbb{A}_{\psi }=\mathbb{A}%
_{\psi }\Rightarrow \mathbb{A}_{\psi }=\left\langle \mathbf{\psi }|A\mathbf{%
\psi }\right\rangle
\end{equation}%
The relationship between vectors represented on these two basis are given by:%
\begin{equation}
\left\vert v\right\rangle =\left\vert \mathbf{\psi }\right\rangle \mathbf{v}%
_{\psi }=\left\vert \mathbf{\varphi }\right\rangle \mathbb{V}_{\varphi }%
\mathbf{v}_{\psi }=\left\vert \mathbf{\varphi }\right\rangle \mathbf{v}%
_{\varphi }
\end{equation}%
thus 
\begin{equation}
\mathbf{v}_{\varphi }=\mathbb{V}_{\varphi }\mathbf{v}_{\psi
}\Longleftrightarrow \mathbf{v}_{\psi }=\mathbb{V}_{\varphi }^{-1}\mathbf{v}%
_{\varphi }
\end{equation}%
and the relation between the \textquotedblright matrix\textquotedblright\ of 
$\left\vert v\right\rangle $ and the representation of $\left\vert
v\right\rangle $ is given by%
\begin{equation}
\left\langle \mathbf{\varphi }|v\right\rangle =\left\langle \mathbf{\varphi
|\varphi }\right\rangle \mathbf{v}_{\varphi }\Longleftrightarrow \mathbf{v}%
_{\varphi }=\left\langle \mathbf{\varphi |\varphi }\right\rangle
^{-1}\left\langle \mathbf{\varphi }|v\right\rangle
\end{equation}%
and 
\begin{eqnarray}
\mathbf{v}_{\varphi }^{\dagger }\left\langle \mathbf{\varphi |\varphi }%
\right\rangle \mathbf{v}_{\varphi } &=&\mathbb{I}_{r}  \notag \\
\left\langle \mathbf{\varphi }|v\right\rangle ^{\dagger }\left\langle 
\mathbf{\varphi |\varphi }\right\rangle ^{-1}\left\langle \mathbf{\varphi }%
|v\right\rangle &=&\mathbb{I}_{r}
\end{eqnarray}%
Obviously in the orthonormal basis 
\begin{equation}
\mathbf{v}_{\psi }=\left\langle \mathbf{\psi }|v\right\rangle
\end{equation}%
In order to study algebraic properties of operators in the two
representations we introduce the resolutions of the identity 
\begin{eqnarray}
I &=&\left\vert \mathbf{\varphi }\right\rangle \left\langle \mathbf{\varphi }%
|\mathbf{\varphi }\right\rangle ^{-1}\left\langle \mathbf{\varphi }%
\right\vert  \notag \\
I &=&\left\vert \mathbf{\psi }\right\rangle \left\langle \mathbf{\psi }|%
\mathbf{\psi }\right\rangle ^{-1}\left\langle \mathbf{\psi }\right\vert
=\left\vert \mathbf{\psi }\right\rangle \left\langle \mathbf{\psi }%
\right\vert .
\end{eqnarray}%
which gives%
\begin{eqnarray}
A &=&\left\vert \mathbf{\varphi }\right\rangle \left\langle \mathbf{\varphi }%
|\mathbf{\varphi }\right\rangle ^{-1}\left\langle \mathbf{\varphi |}A\mathbf{%
\varphi }\right\rangle \left\langle \mathbf{\varphi }|\mathbf{\varphi }%
\right\rangle ^{-1}\left\langle \mathbf{\varphi }\right\vert  \notag \\
&\equiv &\sum \left\vert \varphi _{k}\right\rangle \left\langle \varphi
_{l}\right\vert \left[ \left\langle \mathbf{\varphi }|\mathbf{\varphi }%
\right\rangle ^{-1}\right] _{kk_{1}}\left[ \left\langle \mathbf{\varphi |}A%
\mathbf{\varphi }\right\rangle \right] _{k_{1}k_{2}}\left[ \left\langle 
\mathbf{\varphi }|\mathbf{\varphi }\right\rangle ^{-1}\right] _{k_{2}l}
\end{eqnarray}%
and 
\begin{equation}
A\left\vert \mathbf{\varphi }\right\rangle =\left\vert \mathbf{\varphi }%
\right\rangle \left\langle \mathbf{\varphi }|\mathbf{\varphi }\right\rangle
^{-1}\left\langle \mathbf{\varphi |}A\mathbf{\varphi }\right\rangle
=\left\vert \mathbf{\varphi }\right\rangle \mathbb{A}_{\varphi }
\end{equation}%
Hence 
\begin{equation}
A=\sum \left\vert \varphi _{k}\right\rangle \left\langle \varphi
_{l}\right\vert \left[ \mathbb{A}_{\varphi }\left\langle \mathbf{\varphi }|%
\mathbf{\varphi }\right\rangle ^{-1}\right] _{kl}
\end{equation}%
Using the resolution of the identity the eigenvalue equation 
\begin{equation}
A\left\vert v_{j}\right\rangle =\lambda _{j}\left\vert v_{j}\right\rangle
\end{equation}%
then becomes 
\begin{equation}
\left\vert \mathbf{\varphi }\right\rangle \left\langle \mathbf{\varphi }|%
\mathbf{\varphi }\right\rangle ^{-1}\left\langle \mathbf{\varphi |}A\mathbf{%
\varphi }\right\rangle \left\langle \mathbf{\varphi }|\mathbf{\varphi }%
\right\rangle ^{-1}\left\langle \mathbf{\varphi }|v_{j}\right\rangle
=\lambda _{j}\left\vert v_{j}\right\rangle
\end{equation}%
and hence 
\begin{eqnarray}
&&\left\langle \mathbf{\varphi }|\mathbf{\varphi }\right\rangle \left\langle 
\mathbf{\varphi }|\mathbf{\varphi }\right\rangle ^{-1}\left\langle \mathbf{%
\varphi |}A\mathbf{\varphi }\right\rangle \left\langle \mathbf{\varphi }|%
\mathbf{\varphi }\right\rangle ^{-1}\left\langle \mathbf{\varphi }%
|v_{j}\right\rangle =\lambda _{j}\left\langle \mathbf{\varphi }%
|v_{j}\right\rangle  \notag \\
&\Rightarrow &\left\langle \mathbf{\varphi |}A\mathbf{\varphi }\right\rangle
\left\langle \mathbf{\varphi }|\mathbf{\varphi }\right\rangle
^{-1}\left\langle \mathbf{\varphi }|v_{j}\right\rangle =\lambda
_{j}\left\langle \mathbf{\varphi }|v_{j}\right\rangle  \notag \\
&\Rightarrow &\left\langle \mathbf{\varphi }|\mathbf{\varphi }\right\rangle
^{-1}\left\langle \mathbf{\varphi |}A\mathbf{\varphi }\right\rangle
\left\langle \mathbf{\varphi }|\mathbf{\varphi }\right\rangle
^{-1}\left\langle \mathbf{\varphi }|v_{j}\right\rangle =\lambda
_{j}\left\langle \mathbf{\varphi }|\mathbf{\varphi }\right\rangle
^{-1}\left\langle \mathbf{\varphi }|v_{j}\right\rangle
\end{eqnarray}%
which we can also write as expected as 
\begin{equation}
\mathbb{A}_{\varphi }\mathbf{v}_{\varphi j}=\lambda _{j}\mathbf{v}_{\varphi
j}
\end{equation}%
giving 
\begin{equation}
\mathbb{A}_{\varphi }\mathbf{V}_{\varphi }=\mathbf{V}_{\varphi }\mathbf{%
\lambda }
\end{equation}%
Solutions to this problem can also be obtained by solving the matrix
eigenvalue problem 
\begin{equation}
\left\langle \mathbf{\varphi |}A\mathbf{\varphi }\right\rangle \mathbf{a}%
_{j}=\lambda _{j}\left\langle \mathbf{\varphi }|\mathbf{\varphi }%
\right\rangle \mathbf{a}_{j}
\end{equation}%
which gives 
\begin{equation}
\mathbf{a}_{j}=\left\langle \mathbf{\varphi }|\mathbf{\varphi }\right\rangle
^{-1}\left\langle \mathbf{\varphi }|v_{j}\right\rangle =\mathbf{v}_{\varphi }
\end{equation}%
and 
\begin{equation}
\mathbf{V}_{\varphi }^{\dagger }\left\langle \mathbf{\varphi |}A\mathbf{%
\varphi }\right\rangle \mathbf{V}_{\varphi }=\mathbf{\lambda }
\end{equation}

\subsubsection{Symmetric orthogonalization}

We can choose $\mathbb{V}_{\varphi }$ to be given by 
\begin{equation}
\mathbb{V}_{\varphi }=\mathbb{S}^{-\frac{1}{2}}
\end{equation}%
then 
\begin{equation}
\left\langle \mathbf{\psi }|\mathbf{\psi }\right\rangle =\mathbb{S}^{-\frac{1%
}{2}}\left\langle \mathbf{\varphi }|\mathbf{\varphi }\right\rangle \mathbb{S}%
^{-\frac{1}{2}}=\mathbb{I}_{r}
\end{equation}%
and the eigenvalue problem in this orthonormal representation becomes 
\begin{equation}
\mathbb{S}^{\frac{1}{2}}\mathbb{A}_{\varphi }\mathbb{S}^{-\frac{1}{2}}%
\mathbf{v}_{\psi j}=\lambda _{j}\mathbf{v}_{\psi j}
\end{equation}%
giving 
\begin{equation}
\mathbf{V}_{\psi }^{\dagger }\mathbb{S}^{\frac{1}{2}}\mathbb{A}_{\varphi }%
\mathbb{S}^{-\frac{1}{2}}\mathbf{V}_{\psi }=\mathbf{\lambda }
\end{equation}%
where 
\begin{equation}
\mathbf{V}_{\psi }^{\dagger }\mathbf{V}_{\psi }=\mathbf{V}_{\psi }\mathbf{V}%
_{\psi }^{\dagger }=\mathbb{I}_{r}
\end{equation}

\end{document}
